% !TeX root = bash.tex
\section{Intro}
\begin{frame}
\frametitle{What is shell?}
\begin{figure}[h]
    \centering
    \includegraphics[height=2.5cm]{shell.png}
\end{figure}
\begin{itemize}
    \item A \textbf{command-line interpreter} that reads commands you type and executes them
    \item Also a type of \textbf{scripting language} for automation
    \item The default shell on most Linux distributions is \textbf{bash}(what we are gonna talk about today)
\end{itemize}
\end{frame}

\begin{frame}
\frametitle{The calling stack}
\begin{columns}
\begin{column}{0.45\textwidth}
    \begin{figure}[h]
        \centering
        \includegraphics[width=\textwidth]{bash_kernel_relate}
    \end{figure}
\end{column}
\begin{column}{0.55\textwidth}
    \begin{itemize}
        \item \textbf{User}: You! The person typing commands
        \item \textbf{Terminal}: A program that displays text and captures your input
            (e.g., GNOME Terminal, Windows Terminal, kitty)
        \item \textbf{Bash}: The shell that interprets your commands
        \item \textbf{Kernel}: The core of the OS that actually runs programs,
            manages hardware, files, processes, etc.
    \end{itemize}
\end{column}
\end{columns}
\end{frame}

\begin{frame}
\frametitle{A brief history of bash}
\begin{figure}[h]
    \centering
    \includegraphics[height=3cm]{bash_logo}
\end{figure}
\begin{itemize}
    \item Born: 1989
    \item Runs on Unix-like environments
\end{itemize}
\end{frame}

\begin{frame}
\frametitle{A brief history of Unix}
\begin{figure}[h]
    \centering
    \includegraphics[height=3cm]{unix_logo}
\end{figure}
\begin{itemize}
    \item Born: 1969
    \item Gave rise to Linux, BSD, and Mac OS
    \item We call them ``Unix-like''
\end{itemize}
\end{frame}

\begin{frame}
\frametitle{Unix: The Good Part}
The Unix philosophy (paraphrased):
\begin{itemize}
    \item Everything is a file
    \item hierarchy file system
    \item One tool does one thing
    \item Tools together strong
\end{itemize}
\begin{block}{Note}
    Bash is built upon Unix philosophy. So to learn bash, we first need to know how it works.  
\end{block}
\end{frame}

\begin{frame}
\frametitle{Other Popular Shells}
While this workshop focuses on Bash, there are other shells you might encounter:
\begin{itemize}
    \item \textbf{Zsh}: Highly customizable, compatible with Bash, and features powerful 
        tab-completion and theming (e.g., Oh My Zsh).
    \item \textbf{Fish}: A user-friendly shell with modern syntax, auto-suggestions, 
        and a rich set of features out of the box.
    \item \textbf{Dash}: A lightweight POSIX-compliant shell often used for system scripts 
        due to its speed.
\end{itemize}
Knowing about them helps you choose the right tool for your workflow.
\end{frame}

\begin{frame}[fragile]
\frametitle{Get Your hands dirty}
If you are using Windows, you have two main options to follow along:
\begin{itemize}
    \item \textbf{WSL (Windows Subsystem for Linux)}: Provides a full Linux environment. 
        Install it from the Microsoft Store and choose a distribution (e.g., Ubuntu).
    \item \textbf{Git Bash}: A lightweight Bash emulation that comes with Git for Windows.
        It includes many common Unix commands and is sufficient for this workshop.
\end{itemize}
All commands in this workshop work in both environments.

\bigskip

If you are on Linux or macOS but using a different shell (e.g., Zsh, Fish), you can switch to Bash by typing \texttt{bash} and pressing Enter. To return to your original shell, use \texttt{exit}.
\end{frame}

\begin{frame}[fragile]
\frametitle{Before we start}
\begin{itemize}
    \item This workshop focuses on bash. The syntax is different with PowerShell, fish, etc. 
    \item \tt{\$} indicates the following is a \textbf{bash command}.
        Do not type the \tt{\$}.
    \item \tt{\#} indicates a \textbf{comment}.
        Do not type it or anything after it.
\end{itemize}
For example, when you see:
\begin{lstlisting}[language=bash]
$ echo hello  # print "hello" to stdout
\end{lstlisting}
You are going to type:
\begin{lstlisting}[language=bash]
echo hello
\end{lstlisting}
Then hit Enter. 
\end{frame}
